\section*{Chapter \ref{ch:b3:bioinorganic} --- Bioinorganic Chemistry}

\begin{solution}{exo:b3:bioinorganic:1}
$\log(\theta/(1 - \theta))$ goes from $\log 0.25 = -0.602$ to $\log 4 = +0.602$ while $\log p$ goes up by
$\log(5.9/2.1) = 0.449$: $n = 1.204/0.449 = 2.7$.
\end{solution}

\begin{solution}{exo:b3:bioinorganic:2}
$\ce{Mn^2+} < \ce{Fe^2+} < \ce{Co^2+} < \ce{Ni^2+} < \ce{Cu^2+} > \ce{Zn^2+}$. Zinc ($d^{10}$) has no ligand-field stabilisation and no
Jahn--Teller distortion: it falls back below copper.
\end{solution}

\begin{solution}{exo:b3:bioinorganic:3}
$\ce{Ca^2+}$ and $\ce{Fe^3+}$, hard: carboxylates (and phenolates for iron). $\ce{Cu+}$ and $\ce{Hg^2+}$, soft:
cysteine thiolates. $\ce{Zn^2+}$, borderline: histidine and cysteine.
\end{solution}

\begin{solution}{exo:b3:bioinorganic:4}
$\ce{2H2O -> O2 + 4H+ + 4e-}$: four electrons, four photons per $\ce{O2}$; four moles of photons per
mole of $\ce{O2}$.
\end{solution}

\begin{solution}{exo:b3:bioinorganic:5}
Myoglobin: $\theta = 13/13.37 = 0.972$ and $5/5.37 = 0.931$: it gives up $(0.972 - 0.931)/0.972 = 4$\,\% of its
oxygen. Haemoglobin: $\theta = 0.972$ and $0.724$: it gives up 26\,\%.
\end{solution}

\begin{solution}{exo:b3:bioinorganic:6}
$[\mathrm{HbCO}]/[\mathrm{HbO_2}] = 230 \times \num{50e-6}/0.2095 = 0.055$; fraction $0.055/1.055 = 5.2$\,\%.
\end{solution}

\begin{solution}{exo:b3:bioinorganic:7}
Sulfides $-2$ each, cysteinates $-1$ each. $\ce{[Fe2S2(SR)4]^2-}$: irons sum $-2 + 4 + 4 = +6$: two
Fe(III). $\ce{[Fe2S2(SR)4]^3-}$: $+5$, one Fe(III) and one Fe(II) (mixed valence). $\ce{[Fe4S4(SR)4]^2-}$:
$-2 + 8 + 4 = +10$, two Fe(III) and two Fe(II), each iron $+2.5$ on average.
\end{solution}

\begin{solution}{exo:b3:bioinorganic:8}
$\ce{[PtCl2(NH3)2] + H2O -> [PtCl(H2O)(NH3)2]+ + Cl-}$. The high chloride concentration of blood
pushes the equilibrium back; inside cells chloride is much lower, and the aqua
complex forms. Platinum binds N7 of guanine. In the \textit{trans} isomer the two
\omterm{def:b3:complex-mechanisms:labile}{labile} sites are opposite each other and cannot reach two adjacent bases on one
strand.
\end{solution}

\begin{solution}{exo:b3:bioinorganic:9}
$1/T_1 = 1/\qty{1.2}{s} + \qty{4.0}{L.mmol^{-1}.s^{-1}} \times \qty{0.10}{mmol/L} = 0.833 + 0.400 = \qty{1.23}{s^{-1}}$: $T_1 = \qty{0.81}{s}$.
\end{solution}

\begin{solution}{exo:b3:bioinorganic:10}
$K = K_{\mathrm{PbL}}/K_{\mathrm{CaL}} = 10^{18.0 - 10.7} = 10^{7.3} = \num{2e7}$: lead displaces calcium completely. The free
ligand would also bind the calcium of the blood and lower its concentration
dangerously; the calcium complex exchanges with lead and leaves the calcium
level unchanged.
\end{solution}

\begin{solution}{exo:b3:bioinorganic:11}
Far below $K_M$ the reaction is first order: $k' = (k_{\mathrm{cat}}/K_M)[\mathrm E]$, $t_{1/2} = \ln 2/k'$. Carbonic
anhydrase: $k' = \qty{100}{s^{-1}}$, $t_{1/2} = \qty{6.9}{ms}$. Typical \omterm{def:b3:complex-kinetics:enzyme}{enzyme}: $k' = \qty{0.10}{s^{-1}}$, $t_{1/2} = \qty{6.9}{s}$, a thousand
times longer, too slow for the second a red cell spends in a lung capillary.
\end{solution}

\begin{solution}{exo:b3:bioinorganic:12}
Half the sites with CO means $[\mathrm{HbCO}] = [\mathrm{HbO_2}]$: $p(\ce{CO}) = p(\ce{O2})/M = 0.2095/230 = \qty{9.1e-4}{bar}$, about
\qty{910}{ppm}.
\end{solution}

\begin{solution}{pb:b3:bioinorganic:1}
\textbf{1.} $\theta = 13^{2.7}/(3.5^{2.7} + 13^{2.7}) = 0.972$.

\textbf{2.} $\theta(\qty{5.0}{kPa}) = 0.724$.

\textbf{3.} $5.0/(0.37 + 5.0) = 0.931$.

\textbf{4.} Its $p_{50}$ is ten times lower: at the pressures of a working muscle it is far more
saturated than haemoglobin at the same pressure, so oxygen passes from
haemoglobin to myoglobin.

\textbf{5.} $n = 1$.

\textbf{6.} Binding at one \omterm{def:b3:bioinorganic:haem}{haem} moves its iron into the plane and pulls the histidine and
its helix; the subunits shift together into a form of higher affinity for the
other sites.

\textbf{7.} $\qty{150}{g/L}/\qty{64500}{g/mol} = \qty{2.33e-3}{mol/L}$.

\textbf{8.} $4 \times \num{2.33e-3} = \qty{9.30e-3}{mol/L}$.

\textbf{9.} $\qty{9.30}{mmol/L} \times 0.972 = \qty{9.04}{mmol}$.

\textbf{10.} $\qty{9.30}{mmol/L} \times (0.972 - 0.724) = \qty{2.31}{mmol}$.

\textbf{11.} $\qty{2.31e-3}{mol} \times \qty{22.41}{L/mol} = \qty{52}{mL}$.

\textbf{12.} $0.248/0.972 = 26$\,\%.

\textbf{13.} $5.0^{2.7}/(4.0^{2.7} + 5.0^{2.7}) = 0.646$.

\textbf{14.} $\qty{9.30}{mmol/L} \times (0.972 - 0.646) = \qty{3.03}{mmol}$.

\textbf{15.} $3.03/2.31 = 1.31$: thirty per cent more.

\textbf{16.} Carbon dioxide made by respiration, hydrated to carbonic acid by
carbonic anhydrase, and lactic acid in hard work.

\textbf{17.} $\qty{3.03}{mmol/L} \times \qty{5.0}{L/min} = \qty{15.1}{mmol/min}$.

\textbf{18.} $\qty{15.1e-3}{mol/min} \times \qty{22.41}{L/mol} = \qty{0.34}{L/min}$.

\textbf{19.} $230 \times \num{100e-6}/0.2095 = 0.110$.

\textbf{20.} $0.110/1.110 = 0.099$: about 10\,\% of the sites.

\textbf{21.} The sites still carrying oxygen, in a molecule that also carries CO, are in
the high-affinity form: the curve shifts to the left and they release less of
their oxygen in the tissues.

\textbf{22.} By the competition relation the ratio $[\mathrm{HbCO}]/[\mathrm{HbO_2}]$ falls in proportion to $p(\ce{O2})$:
breathing pure oxygen, almost five times the pressure in air, displaces carbon
monoxide faster.

\textbf{23.} By hindering the linear binding of CO and hydrogen-bonding $\ce{O2}$, it keeps $M$ in
the hundreds instead of the much larger value of free \omterm{def:b3:bioinorganic:haem}{haem}; otherwise even the
carbon monoxide made in the body would block much of the haemoglobin.

\textbf{24.} At rest the blood delivers about \qty{0.34}{L} of oxygen per minute.
\end{solution}
